\section{Proofs for \texorpdfstring{Section~\ref{sec:agy}}{closure, agency and interaction}}\label{app:agy}

The proofs follow the order of \cref{sec:agy}. Two auxiliary lemmas, a four-cycle criterion for exact potential games and a bound on exponential tilts, are proved where they are first needed. For a matrix $X$, $\mathcal F(X)=\{z^\ast Xz:\lvert z\rvert=1\}$ denotes its numerical range; it contains the eigenvalues of $X$, $\mathcal F(X^\top)=\mathcal F(X)$ for real $X$, and $\Vert(X-c)^{-1}\Vert\le1/\operatorname{dist}(c,\mathcal F(X))$ for $c\notin\mathcal F(X)$, because $\lvert(X-c)z\rvert\ge\lvert z^\ast(X-c)z\rvert$ for unit $z$.

\begin{proof}[Proof of \cref{thm:agy-mori-zwanzig}]
In the coordinates $(y,x_\perp)=\mathcal Zx$ the recursion \eqref{eq:agy-interior} reads $y_{t+1}=A_{yy}y_t+A_{y\perp}x_{\perp,t}+B_yx_{\sB,t}+\mathfrak n^y_t$ and $x_{\perp,t+1}=A_{\perp y}y_t+A_{\perp\perp}x_{\perp,t}+B_\perp x_{\sB,t}+\mathfrak n^\perp_t$. Variation of constants in the second equation gives
\begin{equation*}
x_{\perp,t}=A_{\perp\perp}^tx_{\perp,0}+\sum_{s=0}^{t-1}A_{\perp\perp}^{t-1-s}\bigl(A_{\perp y}y_s+B_\perp x_{\sB,s}+\mathfrak n^\perp_s\bigr),
\end{equation*}
and substituting this in the first equation, with $l=t-s$, gives \eqref{eq:agy-memory}.

(a) In stationarity $\Cov(y_t,x_{\perp,t})=\mathsf P\Sigma Z^\top=0$, and $\mathfrak n^\perp_t=Z\mathfrak n_t$ is independent of $(y_t,x_{\perp,t})$, so $\Sigma_\perp=Z\Sigma Z^\top$ satisfies
\begin{equation*}
\Sigma_\perp-A_{\perp\perp}\Sigma_\perp A_{\perp\perp}^\top=A_{\perp y}\mathsf P\Sigma\mathsf P^\top A_{\perp y}^\top+ZQ^{\mathfrak n}Z^\top\succeq ZQ^{\mathfrak n}Z^\top\succ0,
\end{equation*}
the last because $Z$ has full row rank. For a left eigenvector $f^\ast$ of $A_{\perp\perp}$ with eigenvalue $\varpi$ this gives $(1-\lvert\varpi\rvert^2)f^\ast\Sigma_\perp f\ge f^\ast ZQ^{\mathfrak n}Z^\top f>0$, so $\lvert\varpi\rvert<1$. The noises $\mathfrak n_s$, $s\ge0$, are independent of $x_0$, hence $\Cov(\mathfrak f_{t+1},y_0)=A_{y\perp}A_{\perp\perp}^tZ\Sigma\mathsf P^\top=0$.

(b) The same computation applies to $\dot x_\perp=-M_{\perp y}y-M_{\perp\perp}x_\perp+\text{noise}$, whose solution $x_\perp(t)=e^{-M_{\perp\perp}t}x_\perp(0)-\int_0^te^{-M_{\perp\perp}(t-t')}M_{\perp y}y(t')\mathrm dt'+\text{noise}$ is substituted in $\dot y=-M_{yy}y-M_{y\perp}x_\perp+\text{noise}$.
\end{proof}

\begin{proof}[Proof of \cref{thm:agy-closure-identity}]
(a) The proof of \cref{thm:agy-mori-zwanzig} uses no property of $A$, $B$ or $Q^{\mathfrak n}$ beyond linearity. For a coordinate selection take $\mathsf P=(I\ 0)$ and $Z=(0\ I)$ in the ordering $(R,\bar R)$; then $A_{y\perp}=W_{R\bar R}$, $A_{\perp\perp}=W_{\bar R\bar R}$ and $A_{\perp y}=W_{\bar RR}$.

(b) As $x_\perp$ ranges over $\R^{\dim x-\dim y}$, $\mathcal Z^{-1}(0,x_\perp)$ ranges over $\ker\mathsf P$, and $A_{y\perp}x_\perp=\mathsf PA\mathcal Z^{-1}(0,x_\perp)$. Hence $A_{y\perp}=0$ exactly when $\mathsf PA$ vanishes on $\ker\mathsf P$, which is (i), whatever the complement. Every memory, initial-state and colored-noise term of \eqref{eq:agy-memory} has $A_{y\perp}$ as its left factor, so (ii) implies (iii); conversely, the initial-state term at $t=0$ is $A_{y\perp}x_{\perp,0}$, which vanishes for every initial law only if $A_{y\perp}=0$. The equivalence of (i) and (iv) is proved in (c). For a coordinate selection, (ii) is $W_{R\bar R}=0$ by (a).

(c) Given $x_t$, $y_{t+1}\sim\Ncal(\mathsf PAx_t,\mathsf PQ^{\mathfrak n}\mathsf P^\top)$; given $y_t$, $x_t$ has covariance $\Sigma_{x\mid y}$, so $\Cov(y_{t+1}\mid y_t)=\mathsf PA\Sigma_{x\mid y}A^\top\mathsf P^\top+\mathsf PQ^{\mathfrak n}\mathsf P^\top$. Since $y_t$ is a function of $x_t$, $\mathfrak d_{\mathrm{cl}}=I(y_{t+1};x_t\mid y_t)$ is half the difference of the logarithms of the two conditional determinants, which is \eqref{eq:agy-closure-defect}. The stationary covariance satisfies $\Sigma\succeq Q^{\mathfrak n}\succ0$. Write $\Sigma_{x\mid y}=\Sigma^{1/2}(I-\mathsf P_U)\Sigma^{1/2}$, with $\mathsf P_U$ the orthogonal projector onto $U=\operatorname{range}(\Sigma^{1/2}\mathsf P^\top)$. A vector $f$ lies in $U^\perp$ exactly when $\mathsf P\Sigma^{1/2}f=0$, so the range of $\Sigma_{x\mid y}$ is $\Sigma^{1/2}U^\perp=\ker\mathsf P$. The matrix in the determinant is the identity plus a matrix similar to the positive semidefinite $(\mathsf PQ^{\mathfrak n}\mathsf P^\top)^{-1/2}\mathsf PA\Sigma_{x\mid y}A^\top\mathsf P^\top(\mathsf PQ^{\mathfrak n}\mathsf P^\top)^{-1/2}$, so $\mathfrak d_{\mathrm{cl}}=0$ exactly when $\mathsf PA\Sigma_{x\mid y}A^\top\mathsf P^\top=0$, that is, when $\mathsf PA$ vanishes on the range of $\Sigma_{x\mid y}$, which is $\ker\mathsf P$; this proves (i)$\Leftrightarrow$(iv). By the chain rule and the Markov property of $x$,
\begin{equation*}
I(y_{t+1};y_{0:t-1}\mid y_t)\le I(y_{t+1};y_{0:t-1},x_t\mid y_t)=\mathfrak d_{\mathrm{cl}}+I(y_{t+1};y_{0:t-1}\mid x_t,y_t)=\mathfrak d_{\mathrm{cl}}.
\end{equation*}
This holds for every history length $t$. In the stationary state the left side equals $I(y_1;y_{-t:-1}\mid y_0)$, and lengthening the history by one round adds the nonnegative term $I(y_1;y_{-t-1}\mid y_{-t:0})$, so the left side is nondecreasing in $t$ and its limit $\mathfrak d_{\mathrm M}$ obeys the same bound.

(d) Because the input path is known and independent of $(x_0,\mathfrak n)$, conditioning on it amounts to treating it as a deterministic affine term. For a Gaussian initial law, $(x_{\perp,t},y_{0:t})$ is then jointly Gaussian, and the law of $x_{\perp,t}$ given $(y_{0:t},x_{\sB,0:t})$ is $\Ncal(\hat x_{\perp,t},\mathcal V_t)$. The pair $(y_{t+1},x_{\perp,t+1})$ is an affine function of $(y_t,x_{\perp,t},x_{\sB,t})$ plus noise independent of the past, so its conditional law given the past is Gaussian, with mean $(A_{yy}y_t+A_{y\perp}\hat x_{\perp,t}+B_yx_{\sB,t},A_{\perp y}y_t+A_{\perp\perp}\hat x_{\perp,t}+B_\perp x_{\sB,t})$ and a covariance built from $\mathcal V_t$ and $\mathcal ZQ^{\mathfrak n}\mathcal Z^\top$. Conditioning on the observed $y_{t+1}$ gives $(\hat x_{\perp,t+1},\mathcal V_{t+1})$; this is the Kalman recursion, and $\mathcal V_t$ does not depend on the data. The predictive law of $y_{t+1}$ is therefore a function of $(y_t,\hat x_{\perp,t})$ and of the input, and $\hat x_{\perp,t+1}$ is a function of these and of $y_{t+1}$, so the pair is a Markov state for the law of $y$ given the input path. Since $\hat x_{\perp,0}$ is a function of $y_0$, the process $(y_t,\hat x_{\perp,t})$ generates the same past as $y$, and adjoining the $\dim x-\dim y$ coordinates of $\hat x_{\perp,t}$ to $y$ gives a Markov process, whose memory residual vanishes. This bounds the closure order.
\end{proof}

\begin{proof}[Proof of \cref{thm:agy-closure-average}]
Use the orthonormal coordinates $y=\phi_C^\top x$ and $x_\perp=Zx$, so that $x=y\phi_C+Z^\top x_\perp$ and the noise stays standard. In these coordinates $M$ becomes
\begin{equation*}
\tilde M=\begin{pmatrix}\xout^M&f_{\mathrm{col}}^\top\\f_{\mathrm{row}}&M_{\perp\perp}\end{pmatrix},\qquad M_{\perp\perp}=ZMZ^\top,
\end{equation*}
since $\phi_C^\top MZ^\top=(ZM^\top\phi_C)^\top$. In case (a) $M$ is symmetric, $\xout^M=\xout$, and $f_{\mathrm{row}}=f_{\mathrm{col}}=ZM\phi_C=Z(M-\xout)\phi_C$; its norm is $\sigma_C$, because $(M-\xout)\phi_C=(I-\phi_C\phi_C^\top)M\phi_C$ is orthogonal to $\phi_C$ and $Z^\top Z=I-\phi_C\phi_C^\top$, and its entries are $(\varepsilon^\Pi_i+a^\partial_i-\xout)/\sqrt n$, because the row sums of $\Pi_{CC}$ are $\varepsilon^\Pi_i+a^\partial_i$.

Kernel. By \cref{thm:agy-mori-zwanzig}(b) the kernel is $\Gamma(u)=f_{\mathrm{col}}^\top e^{-M_{\perp\perp}u}f_{\mathrm{row}}$. Let $\varkappa_\ast$ be $\xin^M$ in case (b) and $\xin$ in case (a). Then $ZM_{\mathrm{sym}}Z^\top\succeq\varkappa_\ast I$: in case (b) by the definition of $\xin^M$, and in case (a) because $ZMZ^\top=Z(\diag(\varepsilon^\Pi)+L_C)Z^\top+Z\diag(a^\partial)Z^\top$, where the first term is at least $\xin I$ by the definition of $\xin$ in \cref{def:col-contrast}, the vectors $Z^\top z$ being orthogonal to $\one$ and supported on $C$, and the second is positive semidefinite. For every $z$, $\frac{\mathrm d}{\mathrm dt}\lvert e^{-M_{\perp\perp}t}z\rvert^2=-2(e^{-M_{\perp\perp}t}z)^\top ZM_{\mathrm{sym}}Z^\top(e^{-M_{\perp\perp}t}z)\le-2\varkappa_\ast\lvert e^{-M_{\perp\perp}t}z\rvert^2$, so $\Vert e^{-M_{\perp\perp}u}\Vert\le e^{-\varkappa_\ast u}$ and $\lvert\Gamma(u)\rvert\le\sigma_{\mathrm{col}}\sigma_{\mathrm{row}}e^{-\varkappa_\ast u}$. In case (a), $M_{\perp\perp}$ is symmetric positive definite, so $\Gamma(u)=f_{\mathrm{row}}^\top e^{-M_{\perp\perp}u}f_{\mathrm{row}}\ge0$. This proves (b)(i) and the kernel bound of (a).

Case (a), slow mode. Let $M$ have eigenvalues $\varpi_1\le\varpi_2\le\dots\le\varpi_n$ with orthonormal eigenvectors $\mathsf u_1,\dots,\mathsf u_n$. By Courant--Fischer, $\varpi_2\ge\min\{f^\top Mf:\lvert f\rvert=1,f\perp\phi_C\}=\lambda_{\min}(ZMZ^\top)\ge\xin$, and $\varpi_1\le\phi_C^\top M\phi_C=\xout<\xin\le\varpi_2$, so $\varpi_1$ is simple. Write $\phi_C=(\cos\angle)\mathsf u_1+(\sin\angle)w$, with $\angle$ the angle between $\phi_C$ and $\mathsf u_1$ (the sign of $\mathsf u_1$ chosen so that $\cos\angle\ge0$) and $w$ a unit vector orthogonal to $\mathsf u_1$. Then $(M-\xout)\phi_C=(\cos\angle)(\varpi_1-\xout)\mathsf u_1+(\sin\angle)(M-\xout)w$, the two terms are orthogonal, and $w$ lies in the span of eigenvectors with eigenvalues at least $\varpi_2>\xout$, so $\sigma_C^2\ge\sin^2\angle(\varpi_2-\xout)^2$ and $\sin\angle\le\sigma_C/(\varpi_2-\xout)\le\sigma_C/(\xin-\xout)$. For the tangent, write $\mathsf u_1=(\cos\angle)\phi_C+(\sin\angle)w'$ with $w'$ a unit vector orthogonal to $\phi_C$, and apply $Z$ to $M\mathsf u_1=\varpi_1\mathsf u_1$: since $Z\phi_C=0$ and $w'=Z^\top Zw'$, this reads $(\sin\angle)(M_{\perp\perp}-\varpi_1)Zw'=-(\cos\angle)f_{\mathrm{row}}$, with $M_{\perp\perp}=ZMZ^\top\succeq\xin I$ and $\lvert Zw'\rvert=1$. The left side has norm at least $(\sin\angle)(\xin-\varpi_1)\ge(\sin\angle)(\xin-\xout)$ and the right side has norm $(\cos\angle)\sigma_C$, so $\cos\angle>0$ and $\tan\angle\le\sigma_C/(\xin-\xout)$. The matrix $(M-\varpi_1)(M-\varpi_2)$ is positive semidefinite, its eigenvalues being $(\varpi_l-\varpi_1)(\varpi_l-\varpi_2)\ge0$, and since $\lvert M\phi_C\rvert^2=\xout^2+\sigma_C^2$,
\begin{equation*}
0\le\phi_C^\top(M-\varpi_1)(M-\varpi_2)\phi_C=\sigma_C^2-(\xout-\varpi_1)(\varpi_2-\xout),
\end{equation*}
so $0\le\xout-\varpi_1\le\sigma_C^2/(\varpi_2-\xout)\le\sigma_C^2/(\xin-\xout)$. The angle bound is a scalar case of the sine theorem of \citet{davis1970rotation}; both bounds are proved here directly.

Case (a), Markov error. The stationary covariance of $\mathrm dx=-Mx\mathrm dt+\sqrt2\mathrm d\mathsf B$ solves $M\Sigma+\Sigma M=2I$, so $\Sigma=M^{-1}$, and with $\mathsf c_l=\phi_C^\top\mathsf u_l$,
\begin{equation*}
\Cov(y_{t+u},y_t)=\phi_C^\top e^{-Mu}M^{-1}\phi_C=\sum_l\frac{\mathsf c_l^2}{\varpi_l}e^{-\varpi_lu}.
\end{equation*}
Let $\mathfrak m=\sum_{l\ge2}(\mathsf c_l^2/\varpi_l)/\sum_l(\mathsf c_l^2/\varpi_l)$, the memory weight. The autocorrelation is $(1-\mathfrak m)e^{-\varpi_1u}$ plus $\mathfrak m$ times a convex combination of $e^{-\varpi_lu}$, $l\ge2$, each in $[0,e^{-\varpi_1u}]$, so $\Xi\le\mathfrak m$. Since $\sum_{l\ge2}\mathsf c_l^2=\sin^2\angle$ and $\mathsf c_1^2=\cos^2\angle>0$, $\mathfrak m\le(\sin^2\angle/\varpi_2)/(\cos^2\angle/\varpi_1)=\tan^2\angle\cdot\varpi_1/\varpi_2$. With $\tan\angle\le\sigma_C/(\xin-\xout)$ and $\varpi_1/\varpi_2\le\xout/\xin$, this gives \eqref{eq:agy-markov-sym}, and no condition on $\sigma_C$ was used.

Case (b)(ii). Write $\mathsf g=\xin^M-\xout^M$ and $T_\perp=M_{\perp\perp}-\xout^MI$. For a unit complex vector $z$, $\operatorname{Re}(z^\ast T_\perp z)=z^\ast ZM_{\mathrm{sym}}Z^\top z-\xout^M\ge\mathsf g$, so $\Vert T_\perp^{-1}\Vert\le1/\mathsf g$. A vector $(1,\mathsf p)$ spans a right-invariant line of $\tilde M$ exactly when $T_\perp\mathsf p=\mathsf p(f_{\mathrm{col}}^\top\mathsf p)-f_{\mathrm{row}}$, that is, when $\mathsf p$ is a fixed point of $\mathcal H(\mathsf p)=T_\perp^{-1}(\mathsf p(f_{\mathrm{col}}^\top\mathsf p)-f_{\mathrm{row}})$. On the real ball of radius $\mathsf r_-=2\sigma_{\mathrm{row}}/(\mathsf g+(\mathsf g^2-4\sigma_{\mathrm{row}}\sigma_{\mathrm{col}})^{1/2})\le2\sigma_{\mathrm{row}}/\mathsf g$, the smaller root of $\sigma_{\mathrm{col}}\mathsf r^2-\mathsf g\mathsf r+\sigma_{\mathrm{row}}=0$, one has $\lvert\mathcal H(\mathsf p)\rvert\le(\sigma_{\mathrm{col}}\mathsf r_-^2+\sigma_{\mathrm{row}})/\mathsf g=\mathsf r_-$, and since $\mathsf p(f_{\mathrm{col}}^\top\mathsf p)-\mathsf p'(f_{\mathrm{col}}^\top\mathsf p')=(\mathsf p-\mathsf p')(f_{\mathrm{col}}^\top\mathsf p)+\mathsf p'(f_{\mathrm{col}}^\top(\mathsf p-\mathsf p'))$, $\mathcal H$ is a contraction there with constant $2\sigma_{\mathrm{col}}\mathsf r_-/\mathsf g=1-(1-4\sigma_{\mathrm{row}}\sigma_{\mathrm{col}}/\mathsf g^2)^{1/2}<1$. Its fixed point $\mathsf p$ is real, $\varpi_1=\xout^M+f_{\mathrm{col}}^\top\mathsf p$ is a real eigenvalue with $\lvert\varpi_1-\xout^M\rvert\le\sigma_{\mathrm{col}}\mathsf r_-\le2\sigma_{\mathrm{row}}\sigma_{\mathrm{col}}/\mathsf g$, and the right eigenvector $\phi_C+Z^\top\mathsf p$ makes with $\phi_C$ an angle of tangent $\lvert\mathsf p\rvert\le2\sigma_{\mathrm{row}}/\mathsf g$. With $\mathcal L_{\mathsf p}=\begin{pmatrix}1&0\\\mathsf p&I\end{pmatrix}$,
\begin{equation*}
\mathcal L_{\mathsf p}^{-1}\tilde M\mathcal L_{\mathsf p}=\begin{pmatrix}\varpi_1&f_{\mathrm{col}}^\top\\0&M_2\end{pmatrix},\qquad M_2=M_{\perp\perp}-\mathsf pf_{\mathrm{col}}^\top,
\end{equation*}
and $\operatorname{Re}\mathcal F(M_2)\ge\xin^M-\lvert\mathsf p\rvert\sigma_{\mathrm{col}}\ge\xin^M-2\sigma_{\mathrm{row}}\sigma_{\mathrm{col}}/\mathsf g$, which bounds the real parts of the other eigenvalues. Since $4\sigma_{\mathrm{row}}\sigma_{\mathrm{col}}<\mathsf g^2$, $\varpi_1\le\xout^M+2\sigma_{\mathrm{row}}\sigma_{\mathrm{col}}/\mathsf g<\xin^M-2\sigma_{\mathrm{row}}\sigma_{\mathrm{col}}/\mathsf g$, so $\varpi_1$ is not an eigenvalue of $M_2$; it is simple and of smallest real part. The same argument for $\tilde M^\top$, with $f_{\mathrm{row}}$ and $f_{\mathrm{col}}$ exchanged and $M_{\perp\perp}^\top$, which has the same symmetric part, gives a left eigenvector at angle of tangent at most $2\sigma_{\mathrm{col}}/\mathsf g$ whose eigenvalue lies within $2\sigma_{\mathrm{row}}\sigma_{\mathrm{col}}/\mathsf g$ of $\xout^M$ and is therefore $\varpi_1$. This is the invariant-subspace argument of \citet{stewart1973error} in the form needed here.

Case (b)(iii). Let $\mathsf k=(M_2^\top-\varpi_1)^{-1}f_{\mathrm{col}}$ and $\mathcal U_{\mathsf k}=\begin{pmatrix}1&\mathsf k^\top\\0&I\end{pmatrix}$. Then $(\mathcal L_{\mathsf p}\mathcal U_{\mathsf k})^{-1}\tilde M(\mathcal L_{\mathsf p}\mathcal U_{\mathsf k})=\diag(\varpi_1,M_2)$, because $\varpi_1\mathsf k^\top+f_{\mathrm{col}}^\top-\mathsf k^\top M_2=0$, and
\begin{equation*}
(\mathcal L_{\mathsf p}\mathcal U_{\mathsf k})^{-1}=\begin{pmatrix}1+\mathsf k^\top\mathsf p&-\mathsf k^\top\\-\mathsf p&I\end{pmatrix}=:\begin{pmatrix}\mathbf a^\top\\\mathbf G\end{pmatrix}.
\end{equation*}
In the coordinates $(\tilde y,\tilde x)=(\mathcal L_{\mathsf p}\mathcal U_{\mathsf k})^{-1}(y,x_\perp)$ the drift is block diagonal, the noise covariance is $2(\mathcal L_{\mathsf p}\mathcal U_{\mathsf k})^{-1}(\mathcal L_{\mathsf p}\mathcal U_{\mathsf k})^{-\top}$, and $y=\tilde y+\mathsf k^\top\tilde x$. The real parts of the eigenvalues of $M$ are positive, because $M_{\mathrm{sym}}\succ0$, so $\varpi_1>0$. The stationary moments are $\Var(\tilde y)=\lvert\mathbf a\rvert^2/\varpi_1$, $\mathbf X=\Cov(\tilde x,\tilde y)=2(M_2+\varpi_1)^{-1}\mathbf G\mathbf a$ and $\mathbf V=\Var(\tilde x)$ with $M_2\mathbf V+\mathbf VM_2^\top=2\mathbf G\mathbf G^\top$. Since $\E[\tilde y_{t+u}\mid\text{past}]=e^{-\varpi_1u}\tilde y_t$ and $\E[\tilde x_{t+u}\mid\text{past}]=e^{-M_2u}\tilde x_t$,
\begin{equation*}
\Corr(y_t,y_{t+u})=e^{-\varpi_1u}\frac{\Cov(\tilde y,y)}{\Var(y)}+\mathfrak r(u),\qquad\mathfrak r(u)=\frac{\mathsf k^\top e^{-M_2u}\Cov(\tilde x,y)}{\Var(y)},
\end{equation*}
and at $u=0$ the two terms add to one, so $\Corr(y_t,y_{t+u})-e^{-\varpi_1u}=-\mathfrak r(0)e^{-\varpi_1u}+\mathfrak r(u)$. The numerical range of $M_2$ lies in the right half plane, so $\Vert e^{-M_2u}\Vert\le1$, and $\Xi\le2\lvert\mathsf k\rvert\lvert\mathbf X+\mathbf V\mathsf k\rvert/\Var(y)$ because $\Cov(\tilde x,y)=\mathbf X+\mathbf V\mathsf k$. Under $\sigma_{\mathrm{row}},\sigma_{\mathrm{col}}\le\mathsf g/4$ the following hold in turn: $\lvert\mathsf p\rvert\le2\sigma_{\mathrm{row}}/\mathsf g\le\frac12$; $\lvert\varpi_1-\xout^M\rvert\le\mathsf g/8$; $\operatorname{Re}\mathcal F(M_2)\ge\xin^M-\mathsf g/8\ge\frac78\xin^M$, using $\xout^M=\phi_C^\top M_{\mathrm{sym}}\phi_C>0$; $\operatorname{dist}(\varpi_1,\mathcal F(M_2^\top))=\operatorname{dist}(\varpi_1,\mathcal F(M_2))\ge\frac34\mathsf g$, so $\lvert\mathsf k\rvert\le4\sigma_{\mathrm{col}}/(3\mathsf g)\le\frac13$ and $\lvert\mathsf k\rvert\lvert\mathsf p\rvert\le\frac16$; evaluating the Lyapunov equation for $\mathbf V$ on its top eigenvector, $\Vert\mathbf V\Vert\le\Vert\mathbf G\Vert^2/(\frac78\xin^M)=(1+\lvert\mathsf p\rvert^2)/(\frac78\xin^M)\le\frac{10}{7\xin^M}$; $\lvert\mathbf G\mathbf a\rvert=\lvert\mathsf p(1+\mathsf k^\top\mathsf p)+\mathsf k\rvert\le(\frac73\sigma_{\mathrm{row}}+\frac43\sigma_{\mathrm{col}})/\mathsf g$; and $\lvert\mathbf X\rvert\le\frac{16}{7}\lvert\mathbf G\mathbf a\rvert/\xin^M$, since $\operatorname{dist}(-\varpi_1,\mathcal F(M_2))\ge\frac78\xin^M$. Hence
\begin{equation*}
\lvert\mathbf X+\mathbf V\mathsf k\rvert\le\Bigl(\frac{16}3\sigma_{\mathrm{row}}+\frac{104}{21}\sigma_{\mathrm{col}}\Bigr)\frac1{\mathsf g\xin^M}\le\frac{16}3\cdot\frac{\sigma_{\mathrm{row}}+\sigma_{\mathrm{col}}}{\mathsf g\xin^M}.
\end{equation*}
For the variance: $\Var(\tilde y)\ge(1-\lvert\mathsf k\rvert\lvert\mathsf p\rvert)^2/\varpi_1\ge(5/6)^2/\varpi_1$, and $\Var(\mathsf k^\top\tilde x)\le\lvert\mathsf k\rvert^2\Vert\mathbf V\Vert\le\frac{10}{63}(\xin^M)^{-1}\le\frac{10}{63}\varpi_1^{-1}$, because $\varpi_1\le\xout^M+\mathsf g/8<\xin^M$; by the triangle inequality in $L^2$, $\Var(y)\ge(\frac56-(\frac{10}{63})^{1/2})^2/\varpi_1\ge0.189/\varpi_1$. Altogether
\begin{equation*}
\Xi\le2\cdot\frac{4\sigma_{\mathrm{col}}}{3\mathsf g}\cdot\frac{16}3\cdot\frac{\sigma_{\mathrm{row}}+\sigma_{\mathrm{col}}}{\mathsf g\xin^M}\cdot\frac{\varpi_1}{0.189}\le\frac{76\sigma_{\mathrm{col}}(\sigma_{\mathrm{col}}+\sigma_{\mathrm{row}})\varpi_1}{\mathsf g^2\xin^M},
\end{equation*}
the numerical factor being $128/(9\cdot0.189)<75.3$. If in addition $\sigma_{\mathrm{row}}\sigma_{\mathrm{col}}\le\xout^M\mathsf g/2$, then $\varpi_1\le\xout^M+2\sigma_{\mathrm{row}}\sigma_{\mathrm{col}}/\mathsf g\le2\xout^M$, which gives the second bound.

Case (b)(iv). $M\one=M_{\mathrm{sym}}\one+M_{\mathrm{anti}}\one$ and $M^\top\one=M_{\mathrm{sym}}\one-M_{\mathrm{anti}}\one$, and $M_{\mathrm{anti}}\one\perp\one$ because $\one^\top M_{\mathrm{anti}}\one=0$; hence $\sigma_{\mathrm{row}},\sigma_{\mathrm{col}}\le\sigma_{\mathrm{sym}}+\lvert M_{\mathrm{anti}}\phi_C\rvert$, and $\lvert M_{\mathrm{anti}}\phi_C\rvert\le\Vert M_{\mathrm{anti}}\Vert=\frac12\Vert M-M^\top\Vert$. For $M_{ii}=\varepsilon_i+\sum_ja_{ij}$ and $M_{ij}=-a_{ij}$, $(M\one)_i=\varepsilon_i$ and $(M^\top\one)_i=\varepsilon_i+\sum_ja_{ij}-\sum_ja_{ji}$, so $2M_{\mathrm{anti}}\one$ is the vector of differences between attention received and attention given, which sums to zero.
\end{proof}

\begin{proof}[Proof of \cref{prop:agy-drift-precision}]
The Lyapunov equation of the stationary law is $M\Sigma+\Sigma M^\top=2I$, so $M\Sigma=I+J$ and $M=(I+J)\Pi$, with $J$ antisymmetric by construction. Then $M_{\mathrm{sym}}=\frac12[(I+J)\Pi+\Pi(I-J)]=\Pi+\frac12(J\Pi-\Pi J)$. For a unit vector $f$, $f^\top\Pi Jf=f^\top J^\top\Pi f=-f^\top J\Pi f$, so $f^\top M_{\mathrm{sym}}f-f^\top\Pi f=f^\top J\Pi f=(J^\top f)^\top\Pi f$. Since $f^\top Jf=0$, the vector $J^\top f$ is orthogonal to $f$, so $(J^\top f)^\top\Pi f=(J^\top f)^\top(I-ff^\top)\Pi f$, whose modulus is at most $\lvert J^\top f\rvert\lvert(I-ff^\top)\Pi f\rvert\le\Vert J\Vert\lvert(I-ff^\top)\Pi f\rvert\le\Vert J\Vert\lvert\Pi f\rvert$. Two functions on the unit sphere of a subspace that differ everywhere by at most a given amount have minima that differ by at most that amount.
\end{proof}

\begin{proof}[Proof of \cref{prop:agy-pooling}]
(a) The likelihoods are conditionally independent given $\eta_C$, so the union posterior is proportional to $p_\circ\prod_ip(s_i\mid\eta_C)$, and Bayes's rule for each member gives $p(s_i\mid\eta_C)\propto q_i/p_\circ$. In natural parameters $P_i^{-1}=P_\circ^{-1}+\Lambda^s_i$, so $\sum_iP_i^{-1}-(n-1)P_\circ^{-1}=P_\circ^{-1}+\sum_i\Lambda^s_i$, and the information vectors combine in the same way. (b) The pool with weights $1/n$ has natural parameters $\frac1n\sum_i(P_\circ^{-1}+\Lambda^s_i)$, and the product of posteriors $\sum_i(P_\circ^{-1}+\Lambda^s_i)$.
\end{proof}

\begin{proof}[Proof of \cref{prop:agy-consensus}]
(a) Since $\mathcal W^{n_{\mathrm c}}$ is entrywise nonnegative and the $\Lambda^s_j$ are positive semidefinite, $\Lambda_i\succeq\Lambda_\circ$. The union posterior has $h^\ast=\Lambda_\circ m_\circ+n\bar h^s$ and $\Lambda^\ast=\Lambda_\circ+n\bar\Lambda^s$, and $\mathcal W^{n_{\mathrm c}}\one=\one$, so $h_i-h^\ast=n\sum_j(\mathcal W^{n_{\mathrm c}})_{ij}(h^s_j-\bar h^s)$. The stacked vector $(h^s_j-\bar h^s)_j$ is orthogonal to the consensus subspace, on which the symmetric matrix $\mathcal W^{n_{\mathrm c}}$ acts as $(\mathcal W-\frac1n\one\one^\top)^{n_{\mathrm c}}$, with norm $\rho_\perp^{n_{\mathrm c}}$; this gives \eqref{eq:agy-consensus-static}, and the precisions are the same computation. For the means, $\Lambda_i(m_i-m^\ast)=(h_i-h^\ast)-(\Lambda_i-\Lambda^\ast)m^\ast$ and $\Vert\Lambda_i^{-1}\Vert\le\Vert\Lambda_\circ^{-1}\Vert$. For the relative entropy, with $x_k$ the eigenvalues of $X_i$, the covariance part of $\KL(q_i\Vert p^\ast)$ is $\frac12\sum_k\mathsf b(x_k)$ with $\mathsf b(x)=(1+x)^{-1}-1+\ln(1+x)$; $\mathsf b(0)=0$ and $\lvert\mathsf b'(x)\rvert=\lvert x\rvert/(1+x)^2\le4\lvert x\rvert$ for $x\ge-\frac12$, so $0\le\mathsf b(x)\le2x^2$, and the mean part is $\frac12(m_i-m^\ast)^\top\Lambda^\ast(m_i-m^\ast)$.

(b) In information form the steady-state centralized filter is $\hat\eta^{\mathrm c}_{t\mid t}=F\hat\eta^{\mathrm c}_{t-1\mid t-1}+\bar P(\sum_jh^s_j(t)-\Lambda^sF\hat\eta^{\mathrm c}_{t-1\mid t-1})$, because its gain is $\bar P(H^s)^\top(R^s)^{-1}$ for the stacked sensing, and member $i$'s update is the same with $n\tilde h_i(t)$ in place of $\sum_jh^s_j(t)=n\bar h^s(t)$. Subtracting,
\begin{equation*}
\mathrm{dev}_i(t)=(I-\bar P\Lambda^s)F\mathrm{dev}_i(t-1)+n\bar P\bigl(\tilde h_i(t)-\bar h^s(t)\bigr),\qquad I-\bar P\Lambda^s=\bar P(\bar P^{-1}-\Lambda^s)=\bar P(\bar P^-)^{-1},
\end{equation*}
so $(I-\bar P\Lambda^s)F=A_{\mathrm{cl}}$. From $\mathrm{dev}(0)=0$, $\mathrm{dev}_i(t)=\sum_{l=0}^{t-1}A_{\mathrm{cl}}^ln\bar P(\tilde h_i(t-l)-\bar h^s(t-l))$. The consensus rounds preserve the average and contract deviations from it, $(\sum_i\lvert\tilde h_i-\bar h^s\rvert^2)^{1/2}\le\rho_\perp^{n_{\mathrm c}}(\sum_i\lvert h^s_i-\bar h^s\rvert^2)^{1/2}$, and summing the geometric series gives \eqref{eq:agy-consensus-dynamic}; Minkowski's inequality in $L^2$ gives the mean-square form. With $h^s_i=\Lambda^s_i\eta+(H^s_i)^\top(R^s_i)^{-1}v_i$, the noise terms are independent across members with covariances $\Lambda^s_i$ and independent of $\eta$, which gives the stated dispersion. The members' filters use the centralized covariances, so their beliefs share the covariance $\bar P$, and the relative entropy between two Gaussian laws with one covariance is half the squared Mahalanobis distance of their means.
\end{proof}

\begin{proof}[Proof of \cref{prop:agy-literal-consensus}]
A level-1 receiver that reads $m_j$ over a noiseless channel as an observation of $\eta$ whose error variance is the posted $1/\Lambda_j$ (\cref{def:fw-levels}, with $\hat H_{ij}=1$, $\hat P_{ij}=1/\Lambda_j$ and $R_c=0$) adds $\Lambda_j$ to its precision and $\Lambda_jm_j=h_j$ to its information vector, which is the update. The matrix $I+\mathrm{Adj}$ is symmetric, nonnegative and primitive (connected, with self-loops), so by Perron--Frobenius $(I+\mathrm{Adj})^{n_{\mathrm c}}/\rho(I+\mathrm{Adj})^{n_{\mathrm c}}\to\mathsf v\mathsf v^\top$ as the number of rounds $n_{\mathrm c}\to\infty$, with $\mathsf v>0$; this gives (a) and (b), since a connected graph on $n\ge2$ vertices has minimum degree at least one, so every row sum of $I+\mathrm{Adj}$ is at least $2$ and $\rho(I+\mathrm{Adj})\ge2$. For (c), compare the coefficients of $s_j$ in the two means: $\mathsf v_j\Lambda^s_j/\sum_l\mathsf v_l(\Lambda_\circ+\Lambda^s_l)$ against $\Lambda^s_j/(\Lambda_\circ+\sum_l\Lambda^s_l)$. Equality for every $j$ forces $\mathsf v$ to be constant, hence equal row sums of $I+\mathrm{Adj}$, a regular graph, and then the coefficients are $\Lambda^s_j/(n\Lambda_\circ+\sum_l\Lambda^s_l)$, equal to the union's only if $n\Lambda_\circ=\Lambda_\circ$, that is, $\Lambda_\circ=0$ for $n\ge2$; the converse is immediate, and the prior weights follow in the same way. For (d), the weights $\mathsf w_\circ,\mathsf w_j$ sum to one, so the limit minus $\eta$ is $\mathsf w_\circ(m_\circ-\eta)+\sum_j\mathsf w_jv_j$, with independent terms of variances $1/\Lambda_\circ$ and $1/\Lambda^s_j$; the union posterior mean has the least mean-squared error among all estimators, and its error is the posterior variance.
\end{proof}

\begin{proof}[Proof of \cref{lem:agy-decomposition}]
For slot sets $A$ and $A'$ opened by different members, conditional independence of the slots given $\eta_C$ gives $I(o_A;o_{A'}\mid\eta_C)=0$, so the chain rule applied to $I(\eta_C,o_{A'};o_A)$ in two orders gives, under every belief,
\begin{equation}\label{eq:agy-chain}
I(\eta_C;o_{A\cup A'})=I(\eta_C;o_{A'})+I(\eta_C;o_A)-I(o_A;o_{A'}).
\end{equation}
Apply \eqref{eq:agy-chain} to \eqref{eq:agy-member-efe} under $q_i$ and to \eqref{eq:agy-group-efe} under $q_C$, with $A=\mathfrak B_i(\pi_i)$ and $A'=A_{-i}$, and subtract. The terms that involve only $o_{A_{-i}}$, member $i$'s pragmatic terms for the others' slots and the others' controller terms in $G_C$ do not depend on $\pi_i$; what remains is $\mathcal T_i(\pi_i)+\mathcal R_i(\pi)$. Hence $G_i-G_C-\mathcal T_i-\mathcal R_i$ is a function of $\pi_{-i}$, and so is $G_i-\Phi-\mathcal R_i$, which differs from it by $\sum_{j\ne i}\mathcal T_j$.
\end{proof}

\begin{lemma}[Four-cycle criterion, after \citet{monderer1996potential}]\label{lem:agy-four-cycle}
Let $G_i:\mathcal A\to\R$ for $i\in C$. If $\partial^2_{ij}G_i=\partial^2_{ij}G_j$ for all members $i\ne j$, all second differences that start from the null policies ($\pi_i=\pi_j=0$) and all policies of the others, then for any ordering $1,\dots,n$ of the members the function
\begin{equation*}
\Phi'(\pi)=\sum_{k=1}^n\bigl[G_k(\pi^{(k)})-G_k(\pi^{(k-1)})\bigr],\qquad\pi^{(k)}=(\pi_1,\dots,\pi_k,0,\dots,0),
\end{equation*}
is an exact potential: $G_i-\Phi'$ does not depend on $\pi_i$. Conversely, the members of an exact potential game have $\partial^2_{ij}G_i=\partial^2_{ij}G_j$ for all second differences.
\end{lemma}
% src: Monderer and Shapley's four-cycle criterion, in the form used by B4-L2 and L2-P15

\begin{proof}
Let $\mathsf h_k(\pi)=G_k(\pi)-G_k(0,\pi_{-k})$, so that $\Phi'(\pi)=\sum_k\mathsf h_k(\pi^{(k)})$. Fix $i$. The terms with $k<i$ do not involve $\pi_i$. For $k>i$ let $O_k$ be the policies $(\pi_1,\dots,\pi_{i-1},\pi_{i+1},\dots,\pi_{k-1},0,\dots)$ of the members other than $i$ and $k$. The hypothesis for the pair $(i,k)$, with $i$ moving from $0$ to $\pi_i$ and $k$ from $0$ to $\pi_k$ at the context $O_k$, reads
\begin{equation*}
\mathsf h_k(\pi_i,\pi_k,O_k)-\mathsf h_k(0,\pi_k,O_k)=\mathsf h_i(\pi_i,\pi_k,O_k)-\mathsf h_i(\pi_i,0,O_k),
\end{equation*}
and the profiles on the right are $\pi^{(k)}$ and $\pi^{(k-1)}$. Summing over $k>i$, the right sides telescope to $\mathsf h_i(\pi)-\mathsf h_i(\pi^{(i)})$, while each $\mathsf h_k(0,\pi_k,O_k)$ does not involve $\pi_i$. Hence $\Phi'(\pi)-\mathsf h_i(\pi)$ does not depend on $\pi_i$, and neither does $G_i(\pi)-\mathsf h_i(\pi)=G_i(0,\pi_{-i})$. Conversely, if $G_i-\Phi'$ depends only on $\pi_{-i}$ for every $i$, its second difference in $(\pi_i,\pi_j)$ vanishes, so $\partial^2_{ij}G_i=\partial^2_{ij}\Phi'=\partial^2_{ij}G_j$.
\end{proof}

\begin{proof}[Proof of \cref{thm:agy-agency}]
(a) A move of member $i$ from $\pi_i$ to $\pi'_i$ at fixed $\pi_{-i}$ has probability
\begin{equation*}
n^{-1}E_i(\pi'_i)e^{-\gamma G_i(\pi'_i,\pi_{-i})}/\mathsf Z_i(\pi_{-i}),
\end{equation*}
with a normalizer that depends on $\pi_{-i}$ only. Along the closed path, the four normalizers, the four habit factors and the factors $n^{-1}$ appear once in each direction and cancel, and the exponents give the forward-to-backward ratio $\exp(-\gamma(\partial^2_{ij}G_j-\partial^2_{ij}G_i))$. By \cref{lem:agy-decomposition}, $\partial^2_{ij}G_i=\partial^2_{ij}\Phi+\partial^2_{ij}\mathcal R_i$, because a function of $\pi_{-i}$ has zero second difference in $(\pi_i,\pi_j)$, and likewise for $j$; this gives $\mathcal L_{ij}$. Through the null policies $\mathcal R_i(0,\cdot)=0$, so by the chain rule
\begin{equation*}
\partial^2_{ij}\mathcal R_i=\bigl[I_{q_i}-I_{q_C}\bigr]\bigl(o_{A'};o_{B'},o_{A_{-ij}}\bigr)-\bigl[I_{q_i}-I_{q_C}\bigr]\bigl(o_{A'};o_{A_{-ij}}\bigr)=\bigl[I_{q_i}-I_{q_C}\bigr]\bigl(o_{A'};o_{B'}\mid o_{A_{-ij}}\bigr),
\end{equation*}
and the same holds for $j$; subtracting gives \eqref{eq:agy-cycle}, in which the terms in $q_C$ cancel.

(b) If the chain is reversible, multiplying the detailed-balance equations along any closed path shows that it has the same probability in both directions, which is the easy half of Kolmogorov's criterion \citep{kelly1979reversibility}; so every $\mathcal L_{ij}$ through the null policies vanishes, which is agreement on redundancy. If the members agree, then $\partial^2_{ij}G_j-\partial^2_{ij}G_i=\mathcal L_{ij}=0$ for every second difference from the null policies, and \cref{lem:agy-four-cycle} gives an exact potential $\Phi'$. If $\Phi'$ is an exact potential, a move of $i$ from $\pi$ to $\pi'=(\pi'_i,\pi_{-i})$ satisfies detailed balance for $\mu\propto\prod_kE_ke^{-\gamma\Phi'}$: the ratio $\mu(\pi)\mathsf T(\pi,\pi')/(\mu(\pi')\mathsf T(\pi',\pi))$ equals $\exp(-\gamma[\Phi'(\pi)-\Phi'(\pi')+G_i(\pi')-G_i(\pi)])=1$, since $G_i-\Phi'$ does not depend on $\pi_i$. Finally an exact potential game has $\mathcal L_{ij}=0$ by the converse in \cref{lem:agy-four-cycle}.

(c) The chain is irreducible, since single moves with positive probability connect all profiles, and aperiodic, since it has self-loops, so its stationary law is unique and equals $\mu$. In the potential of \cref{lem:agy-four-cycle}, $G_k(\pi^{(k)})-G_k(\pi^{(k-1)})=-I_{q_k}(\eta_C;o_{A_k}\mid o_{A_1},\dots,o_{A_{k-1}})+\psi^{(k)}_{A_k}(q_k)$, because the pragmatic terms of the other slots cancel, and \eqref{eq:agy-chain} turns the conditional information into the first and third terms of \eqref{eq:agy-potential}. Now let $q_C$ have the agreed conditional redundancies. By the chain rule, $I_{q_k}(o_{A_k};o_{A_1},\dots,o_{A_{k-1}})=\sum_{l<k}I_{q_k}(o_{A_k};o_{A_l}\mid o_{A_1},\dots,o_{A_{l-1}})$, and each term is a conditional redundancy between the slots of members $k$ and $l$ in a context opened by other members, so it equals the corresponding term under $q_C$. The charges therefore add up to $\sum_kI_{q_C}(o_{A_k};o_{A_1},\dots,o_{A_{k-1}})$, which by \eqref{eq:agy-chain}, applied $n-1$ times, is $-I_{q_C}(\eta_C;o_{\mathfrak B(\pi)})+\sum_kI_{q_C}(\eta_C;o_{A_k})$. Hence $\Phi'=G_C+\sum_k\mathcal T_k$, with the $\mathcal T_k$ of \cref{lem:agy-decomposition}, and $\mathbb P^\infty\propto\prod_k[E_ke^{-\gamma\mathcal T_k}]e^{-\gamma G_C}$. For two members the only redundancies are $I(o_{A_1};o_{A_2})$ for all policies, on which the members agree, so $q_C=q_1$ has them.

(d) Write $o_A=\mathsf H_A\eta_C+\mathsf n_A$ with noise $\mathsf n_A$ independent of $\eta_C$, so that the law of the noise is the same under every belief. Under $q_i$ the pair $(\eta_C,o)$ is then the image of a pair distributed under $q$ by the translation $\eta_C\mapsto\eta_C+m_i$, $o_A\mapsto o_A+\mathsf H_Am_i$, which acts on each variable separately. Mutual informations and conditional mutual informations are invariant under bijective measurable maps of each argument, so every information term is the same under every $q_i$ and under every translate $q_C$. Then $\mathcal R_i\equiv0$, $q_C$ has the agreed redundancies, the epistemic part of $\mathcal T_i$ vanishes, and (c) gives the group posterior. Under (P0), $q_i=q_C$ for every $i$, so for any slot likelihoods $\mathcal R_i\equiv0$ and $\mathcal T_i\equiv0$, and the chain rule for mutual information writes $-I_{q_C}(\eta_C;o_{\mathfrak B(\pi)})=-I_{q_C}(\eta_C;o_{\mathfrak B_i(\pi_i)}\mid o_{A_{-i}})-I_{q_C}(\eta_C;o_{A_{-i}})$, in which the second term depends on $\pi_{-i}$ only, so the part that depends on $\pi_i$ is the conditional information gain.
\end{proof}

\begin{lemma}[Tilts]\label{lem:agy-tilt}
Let $\mathbb Q$ be a law on a finite set and $g$ a function with oscillation $\operatorname{osc}g$, and let $\mathbb Q_g\propto\mathbb Qe^g$. Then $\mathrm{TV}(\mathbb Q,\mathbb Q_g)\le\tanh(\operatorname{osc}g/4)$, with equality for a law with masses $(1+e^{\operatorname{osc}g/2})^{-1}$ and $e^{\operatorname{osc}g/2}(1+e^{\operatorname{osc}g/2})^{-1}$ on a maximizer and a minimizer of $g$.
\end{lemma}
% src: B4-L3

\begin{proof}
Shift $g$ so that its maximum is $\frac12\operatorname{osc}g$ and its minimum $-\frac12\operatorname{osc}g$, and put $\mathsf u=e^{\operatorname{osc}g/2}$. For a set $S$ with $\mathbb Q(S)=\theta_S$, $\mathbb Q_g(S)$ is largest when $g$ is maximal on $S$ and minimal off $S$, so $\mathbb Q_g(S)-\theta_S\le\theta_S\mathsf u/(\theta_S\mathsf u+(1-\theta_S)\mathsf u^{-1})-\theta_S=\mathsf v\theta_S(1-\theta_S)/(1+\mathsf v\theta_S)$ with $\mathsf v=\mathsf u^2-1$. The derivative in $\theta_S$ vanishes where $\mathsf v\theta_S^2+2\theta_S-1=0$, at $\theta_S=1/(1+\mathsf u)$, where the value is $(\mathsf u-1)/(\mathsf u+1)=\tanh(\operatorname{osc}g/4)$. The two-point law attains it.
\end{proof}

\begin{proof}[Proof of \cref{thm:agy-heterogeneous}]
Ratio. By \cref{lem:agy-decomposition}, member $i$'s conditional at $\pi_{-i}$ is its conditional in $\mathsf T_\Phi$ tilted by $e^{-\gamma\mathcal R_i(\cdot,\pi_{-i})}$, so the two probabilities of moving to $\pi'_i$ differ by the factor
\begin{equation*}
e^{-\gamma\mathcal R_i(\pi'_i,\pi_{-i})}\big/\E e^{-\gamma\mathcal R_i(\cdot,\pi_{-i})},
\end{equation*}
with the expectation under the reference conditional. Numerator and denominator both lie in $[e^{-\gamma\max\mathcal R_i},e^{-\gamma\min\mathcal R_i}]$, so the factor lies in $[e^{-\gamma\sigma_i},e^{\gamma\sigma_i}]$. An off-diagonal entry of $\mathsf T$ is $n^{-1}$ times one such probability, and a diagonal entry is an average of such probabilities, so both inherit the bound.

(a) By the Markov chain tree theorem \citep{anantharam1989proof}, the stationary probability of a profile is proportional to the sum, over spanning trees of the profiles directed toward it, of the products of the transition probabilities along their $\lvert\mathcal A\rvert-1$ edges, all off the diagonal. $\mathsf T$ and $\mathsf T_\Phi$ have the same off-diagonal support, so each tree weight changes by a factor in $[e^{-(\lvert\mathcal A\rvert-1)\gamma\sigma_{\max}},e^{(\lvert\mathcal A\rvert-1)\gamma\sigma_{\max}}]$, and so do the unnormalized stationary weights; after normalization $\ln(\mathbb P^\infty/\mathbb P_\Phi)$ takes values in an interval of length at most $2(\lvert\mathcal A\rvert-1)\gamma\sigma_{\max}$ that contains $0$. So $\mathbb P^\infty$ is a tilt of $\mathbb P_\Phi$ by a function of oscillation at most $2(\lvert\mathcal A\rvert-1)\gamma\sigma_{\max}$, \cref{lem:agy-tilt} gives the bound in total variation, and $\KL(\mathbb P^\infty\Vert\mathbb P_\Phi)\le\max\ln(\mathbb P^\infty/\mathbb P_\Phi)$.

(b) Let $d_{\mathrm H}$ be the Hamming distance between profiles and $\mathrm{Wass}_{\mathrm H}$ the corresponding Wasserstein distance between laws; $\mathrm{TV}\le \mathrm{Wass}_{\mathrm H}$ because $d_{\mathrm H}(\pi,\pi')\ge\mathbf 1[\pi\ne\pi']$. Take two profiles that differ only in member $j$'s policy and run one step of $\mathsf T_\Phi$ from each with the same chosen member $i$. If $i=j$ the two conditionals coincide and the coupled draws agree, giving distance $0$; if $i\ne j$, a maximal coupling of the two conditionals of $i$ disagrees with probability at most $\mathfrak I_{ij}$, giving distance one plus that indicator. The expected distance is at most $1-(1-\mathfrak D_\Phi)/n$. Since $d_{\mathrm H}$ is the path metric of single-member changes, $\mathrm{Wass}_{\mathrm H}(\delta_\pi\mathsf T_\Phi,\delta_{\pi'}\mathsf T_\Phi)\le(1-(1-\mathfrak D_\Phi)/n)d_{\mathrm H}(\pi,\pi')$ for all profiles, and coupling initial laws optimally gives $\mathrm{Wass}_{\mathrm H}(\nu\mathsf T_\Phi,\nu'\mathsf T_\Phi)\le(1-(1-\mathfrak D_\Phi)/n)\mathrm{Wass}_{\mathrm H}(\nu,\nu')$; this is the path-coupling argument of \citet{bubley1997path}. From one profile, a step of $\mathsf T$ and a step of $\mathsf T_\Phi$ with the same chosen member $i$ differ only through $i$'s conditional, which by the ratio argument is a tilt of the reference conditional by a function of oscillation at most $\gamma\sigma_i$, so by \cref{lem:agy-tilt} $\mathrm{Wass}_{\mathrm H}(\nu\mathsf T,\nu\mathsf T_\Phi)\le n^{-1}\sum_i\tanh(\gamma\sigma_i/4)$. Then
\begin{equation*}
\mathrm{Wass}_{\mathrm H}(\mathbb P^\infty,\mathbb P_\Phi)=\mathrm{Wass}_{\mathrm H}(\mathbb P^\infty\mathsf T,\mathbb P_\Phi\mathsf T_\Phi)\le\frac1n\sum_i\tanh\frac{\gamma\sigma_i}4+\Bigl(1-\frac{1-\mathfrak D_\Phi}n\Bigr)\mathrm{Wass}_{\mathrm H}(\mathbb P^\infty,\mathbb P_\Phi),
\end{equation*}
which is \eqref{eq:agy-contraction}. Two conditionals of $i$ under $\mathbb P_\Phi$ at profiles that differ in $j$'s policy, from $\pi_j$ to $\pi'_j$, are tilts of each other by $-\gamma[\Phi(\cdot,\pi_j,\cdot)-\Phi(\cdot,\pi'_j,\cdot)]$, whose oscillation in $\pi_i$ is $\gamma\max\lvert\partial^2_{ij}\Phi\rvert=\gamma\max\lvert\partial^2_{ij}G_C\rvert$, because the tilts $\mathcal T_k$ are separable. A second difference of $G_C$ from the null policies is $I_{q_C}(o_{A'};o_{B'}\mid o_{A_{-ij}})\ge0$ (proof of \cref{prop:agy-performative-sign}), and every second difference is the signed sum $\partial^2_{(\pi_i,\pi'_i),(\pi_j,\pi'_j)}=\partial^2_{(0,\pi'_i),(0,\pi'_j)}-\partial^2_{(0,\pi_i),(0,\pi'_j)}-\partial^2_{(0,\pi'_i),(0,\pi_j)}+\partial^2_{(0,\pi_i),(0,\pi_j)}$ of four of them, two with each sign; with two policies per member every second difference is one of them up to sign. Finally $\tanh x\le x$.

(c) $\mathsf T_\Phi$ is the average over members of the resampling operators, each the conditional expectation given $\pi_{-i}$ under $\mathbb P_\Phi$ and hence an orthogonal projection in $L^2(\mathbb P_\Phi)$; so $I-\mathsf T_\Phi$ is self-adjoint and positive, with smallest eigenvalue $\mathrm{gap}_\Phi$ on functions of mean zero. Write $\mathbb P^\infty=\mathbb P_\Phi(1+f)$ and $\mathsf E=\mathsf T-\mathsf T_\Phi$, and let $\mathsf E^\ast$ be the adjoint in $L^2(\mathbb P_\Phi)$, which represents the action of $\mathsf E$ on densities. Stationarity and the self-adjointness of $\mathsf T_\Phi$ give $f-\mathsf T_\Phi f=\mathsf E^\ast(1+f)$, with $\E_{\mathbb P_\Phi}f=0$, and $f-\mathsf T_\Phi f$ also has mean zero. The entrywise bound $\lvert\mathsf E\rvert\le(e^{\gamma\sigma_{\max}}-1)\mathsf T_\Phi$ and the contraction of $\mathsf T_\Phi$ on $L^2(\mathbb P_\Phi)$ give $\Vert\mathsf E^\ast\Vert=\Vert\mathsf E\Vert\le e^{\gamma\sigma_{\max}}-1$, and $\mathsf E^\ast1=\mathbb P_\Phi\mathsf T/\mathbb P_\Phi-1$. Hence $\mathrm{gap}_\Phi\Vert f\Vert\le\Vert\mathsf E^\ast1\Vert+(e^{\gamma\sigma_{\max}}-1)\Vert f\Vert$, which is \eqref{eq:agy-spectral}, with $\Vert\mathsf E^\ast1\Vert\le e^{\gamma\sigma_{\max}}-1$.

For the group posterior, the same arguments apply to the reference chain of $G_C$, whose resampling laws are tilted into the members' conditionals by $-\gamma(G_i-G_C)$.
\end{proof}

\begin{proof}[Proof of \cref{prop:agy-dispersion}]
The definition of $\mathrm{Th}_i$ gives $e^{-\mathrm{Th}_i}P_C\preceq P_i\preceq e^{\mathrm{Th}_i}P_C$. The posterior covariance after slots $O$, $P_{\mid O}=(P^{-1}+\mathsf H_O^\top\mathsf R_O^{-1}\mathsf H_O)^{-1}$, is monotone in $P$, and for $\mathsf c\ge1$, $(\mathsf cP)^{-1}+\mathsf H_O^\top\mathsf R_O^{-1}\mathsf H_O\succeq\mathsf c^{-1}(P^{-1}+\mathsf H_O^\top\mathsf R_O^{-1}\mathsf H_O)$, so $(\mathsf cP)_{\mid O}\preceq\mathsf cP_{\mid O}$, and the reverse holds for $\mathsf c\le1$. Hence $e^{-\mathrm{Th}_i}P_{C\mid O}\preceq P_{i\mid O}\preceq e^{\mathrm{Th}_i}P_{C\mid O}$. For Gaussian slots, $I_q(\eta_C;o_A\mid o_O)=\frac12\ln\det(I+\mathsf R_A^{-1/2}\mathsf H_AP_{q\mid O}\mathsf H_A^\top\mathsf R_A^{-1/2})$ is monotone in $P_{q\mid O}$, and for a positive semidefinite $N$ of size $\dim o_A$, $\det(I+\mathsf cN)$ lies between $\det(I+N)$ and $\mathsf c^{\dim o_A}\det(I+N)$. So $I_{q_i}(\eta_C;o_A\mid o_O)$ and $I_{q_i}(\eta_C;o_A)$ differ from their values under $q_C$ by at most $\frac12\dim(o_A)\mathrm{Th}_i$, and by \eqref{eq:agy-chain} their difference, the redundancy $I(o_A;o_O)$, differs by at most $\dim(o_A)\mathrm{Th}_i$. The second bound holds because both redundancies are nonnegative. Since $\operatorname{osc}\le2\max\lvert\mathcal R_i\rvert$, the bound on $\sigma_i$ follows. The means do not enter because Gaussian mutual informations depend only on covariances.
\end{proof}

\begin{proof}[Proof of \cref{prop:agy-index}]
(a) The habit-only model is admissible, so the infimum is at most the denominator and $0\le\Agency\le1$. A larger class has a smaller infimum. Relabeling members permutes profiles, which leaves relative entropies unchanged, so the index is invariant when the class is closed under relabeling. With free joint habits, $E_C\propto\mathbb Pe^{\gamma_CG_C}$ reproduces any full-support law $\mathbb P$ at one decision state. (b) The law $\mathbb P^\infty$ is then itself in the class. (c) The first inequality is the definition with the group model of $\mathbb P_\Phi$ admissible. For the second, $\KL(\mathbb P^\infty\Vert E)=\KL(\mathbb P^\infty\Vert\mathbb P_\Phi)+\E_{\mathbb P^\infty}\ln(\mathbb P_\Phi/E)$ and $\ln(\mathbb P_\Phi/E)=-\gamma\Phi-\ln\mathsf Z_\Phi$, so
\begin{equation*}
\E_{\mathbb P^\infty}\ln\frac{\mathbb P_\Phi}E=\KL(\mathbb P_\Phi\Vert E)-\gamma\bigl(\E_{\mathbb P^\infty}\Phi-\E_{\mathbb P_\Phi}\Phi\bigr)\ge\KL(\mathbb P_\Phi\Vert E)-\gamma\operatorname{osc}(\Phi)B_{\mathrm{TV}}=\mathsf c_\Phi.
\end{equation*}
For $\mathsf c_\Phi>0$ the map $x\mapsto x/(x+\mathsf c_\Phi)$ is increasing on $[0,\infty)$. Jensen's inequality gives $\KL(\mathbb P^\infty\Vert\mathbb P_\Phi)\le\ln(1+\Vert\mathbb P^\infty/\mathbb P_\Phi-1\Vert^2_{L^2(\mathbb P_\Phi)})$, so the bounds of \cref{thm:agy-heterogeneous} supply $B_{\mathrm{KL}}$ and $B_{\mathrm{TV}}$.
\end{proof}

\begin{proof}[Proof of \cref{prop:agy-synchronous}]
If the previous policies $(\pi_1,\pi_2)$ have a product law, the new ones depend on $\pi_2$ and on $\pi_1$ respectively, through independent randomness, so they are independent: the map induced on laws sends product laws to product laws. Restricted to them it is a continuous self-map of a product of simplices, which has a fixed point by Brouwer's theorem; a fixed product law is stationary. Each member's conditional is proportional to $E_i(\pi_i)e^{-\gamma G_i}$ with full-support habits and finite expected free energies, so every transition probability of the synchronous chain, a product of two such conditionals, is positive, and the chain has a unique stationary law. Positivity is needed: for the deterministic conditionals that copy the other member's previous policy, the chain exchanges the two policies and the uniform law on the profiles with equal policies is stationary without being a product law. A product law has log odds ratio zero, while by \cref{prop:agy-performative-sign} a group posterior with product habits whose two slots from the null policies carry redundant information has log odds ratio $-\gamma_CI_{q_C}(o_{A'};o_{B'})<0$ when $\gamma_C>0$.
\end{proof}

\begin{proof}[Proof of \cref{thm:agy-joint-acts}]
With intentions shared, member $i$'s outcome vector is $(o^\star,o_{\mathfrak B(\pi^{\mathrm{epi}})})$, and its expected free energy is minus the information gain about $(\eta,\eta')$ minus the expected log preference, the second form of the expected free energy of \cref{sec:fw}. Under its model the joint law of $(\eta,\eta',o^\star,o_{\mathfrak B})$ is Gaussian with a covariance that does not depend on $u$ and a mean affine in the effects of $u$; Gaussian information gains depend only on covariances, and the slots observe $\eta$ before the acts, so the information term and the slots' pragmatic terms depend on $\pi^{\mathrm{epi}}$ only. The predictive law of $o^\star$ is Gaussian with mean $\mathsf H_\star Fm_i+f^{(i)}(u)$ and a covariance independent of $u$, so $-\E\ln\tilde p(o^\star)$ is $\frac12\Vert\mathsf H_\star Fm_i+f^{(i)}(u)-o^\star_i\Vert^2_{\Lambda^\star_i}$ plus a constant, which is \eqref{eq:agy-act-efe}.

(a) For $\mathcal B_{\mathrm{int}}=0$, $G^{\mathrm{act}}_k$ is a quadratic function of $u$ whose part bilinear in $(u_i,u_j)$ is
\begin{equation*}
u_i^\top\mathsf M^{(k)\top}_i\Lambda^\star_k\mathsf M^{(k)}_ju_j,\qquad\text{so that}\qquad\partial^2_{ij}G^{\mathrm{act}}_k=(u'_i-u_i)^\top\mathsf M^{(k)\top}_i\Lambda^\star_k\mathsf M^{(k)}_j(u'_j-u_j).
\end{equation*}
The computation of the cycle ratio in the proof of \cref{thm:agy-agency}(a) applies to any profile of expected free energies, and \cref{lem:agy-four-cycle} gives the criterion. (b) Expanding, $\frac12\Vert\mathsf z_i+y\Vert^2_{\Lambda^\star}=\frac12\Vert\mathsf z_C+y\Vert^2_{\Lambda^\star}+\langle\mathsf z_i-\mathsf z_C,y\rangle_{\Lambda^\star}+\mathrm{const}$ with $y=\sum_j\mathsf M_ju_j$, and $\langle\mathsf z_i-\mathsf z_C,y\rangle_{\Lambda^\star}$ splits into $\langle\mathsf z_i-\mathsf z_C,\mathsf M_iu_i\rangle_{\Lambda^\star}$ and terms in $u_{-i}$, which gives \eqref{eq:agy-act-potential}. If the epistemic component has a potential $G^{\mathrm{epi}}_C+\sum_k\mathcal T_k$, then $G_i-(G^{\mathrm{epi}}_C+\sum_k\mathcal T_k)-\Phi^{\mathrm{act}}$ depends on the others' policies only, and the stationary law is the Gibbs law of \cref{thm:agy-agency}(c) for the potential $G_C+\sum_k\mathcal T_k+\sum_j\langle\mathsf z_j-\mathsf z_C,\mathsf M_ju_j\rangle_{\Lambda^\star}$, the group expected free energy plus own-policy terms, which is the group posterior with habits tilted by these terms. (c) $G^{\mathrm{act}}_i-G^{\mathrm{act}}_j=\langle\mathsf z_i-\mathsf z_j,f(u)\rangle_{\Lambda^\star}+\mathrm{const}$, whose second difference is $\langle\mathsf z_i-\mathsf z_j,\mathsf H_\star\partial^2_{ij}\mathcal B_{\mathrm{int}}\rangle_{\Lambda^\star}$, because the additive part of $f$ has zero second difference; the cycle sum is $\partial^2_{ij}G_j-\partial^2_{ij}G_i$, which is \eqref{eq:agy-acts-cycle}. With a common displacement $\mathsf z_i=\mathsf z_C$ and $\Lambda^\star_C=\Lambda^\star$ every $G^{\mathrm{act}}_i$ equals $G^{\mathrm{act}}_C$ up to a constant. Common targets $o^\star_i$ alone do not give this, because $\mathsf z_i=\mathsf H_\star Fm_i-o^\star_i$ also contains the member's predicted outcome.
\end{proof}

\begin{proof}[Proof of \cref{prop:agy-performative-sign}]
The habits are a product and cancel in the odds ratio, which leaves $-\gamma_C[G_C(\pi'_i,\pi'_j)+G_C(\pi_i,\pi_j)-G_C(\pi'_i,\pi_j)-G_C(\pi_i,\pi'_j)]=-\gamma_C\partial^2_{ij}G_C$. For epistemic slots from the null policies, with $O=A_{-ij}$, the pragmatic terms are separable, and \eqref{eq:agy-chain} applied to $(A',B'\cup O)$ and to $(A',O)$ gives $-I(\eta_C;o_{A'\cup B'\cup O})+I(\eta_C;o_{B'\cup O})=-I(\eta_C;o_{A'})+I(o_{A'};o_{B'\cup O})$ and $-I(\eta_C;o_{A'\cup O})+I(\eta_C;o_O)=-I(\eta_C;o_{A'})+I(o_{A'};o_O)$; subtracting, $\partial^2_{ij}G_C=I_{q_C}(o_{A'};o_{B'\cup O})-I_{q_C}(o_{A'};o_O)=I_{q_C}(o_{A'};o_{B'}\mid o_O)$. For additive acts the epistemic part does not depend on $u$ and the act part is the quadratic of \eqref{eq:agy-act-efe}, whose second difference is the bilinear term.
\end{proof}

\begin{proof}[Proof of \cref{prop:agy-identity}]
(a) For each position, the expected change of $\mu_i$ in $\mathrm dt$ given the present is $[-\varepsilon\mu_i+\beta(z-\mu_i)+g(\mathsf fz-\mu_i)]\mathrm dt$ plus a communication term that sums to zero over the members, the term with $g$ being the expected jump at a replacement, which occurs at rate $g$ and takes $\mu_i$ to $\mathsf fz$ plus mean-zero noise. Averaging over positions gives the first row of \eqref{eq:agy-identity}, and the record's drift $\omega(\bar\mu-z)-\delta z$ gives the second. The noise terms are martingale increments, so $\E[(\bar\mu,z)_{t+u}\mid\text{process up to }t]=e^{u\mathsf A}(\bar\mu,z)_t$. With $\mathsf p_1=\varepsilon+\beta+g$, $\mathsf p_2=\omega+\delta$ and $\mathsf E=\omega(\beta+g\mathsf f)\ge0$, the eigenvalues of $\mathsf A$ are $-\varpi_\pm$ with $\varpi_\pm=\frac12[(\mathsf p_1+\mathsf p_2)\pm((\mathsf p_1-\mathsf p_2)^2+4\mathsf E)^{1/2}]$, real, and positive because $\operatorname{tr}\mathsf A<0$ and $\det\mathsf A>0$; the system is stable. The martingale part of $(\bar\mu,z)$ lies in the first coordinate, and in the stationary law its quadratic variation grows at a constant rate $\mathsf q$ that is at least $g$ times the newcomers' noise variance divided by $n$, because a replacement at one position moves $\bar\mu$ by $(\mathsf fz+\xi-\mu_i)/n$ with $\xi$ the newcomer's noise, independent of the present and of mean zero, and replacements occur at total rate $ng$. The stationary covariance $\mathsf S$ of $(\bar\mu,z)$ therefore solves $\mathsf A\mathsf S+\mathsf S\mathsf A^\top+\diag(\mathsf q,0)=0$ with $\mathsf q>0$, and it is positive definite because $\mathsf A$ is stable and $\omega>0$ carries the noise of the first coordinate into the second. The projection of the state on the left eigenvector of $-\varpi_-$ is therefore a nondegenerate random variable, so $\varpi_-$ is the slowest rate at which the conditional expectation of the pair relaxes, and $\tau_{\mathrm{id}}^{-1}=\varpi_-$. When $\beta+g\mathsf f>0$, a right eigenvector of $-\varpi_-$ is $(\beta+g\mathsf f,\mathsf p_1-\varpi_-)$, whose first component is nonzero, and $\varpi_-<\varpi_+$, so the conditional expectation of $\bar\mu$ alone also relaxes at the rate $\varpi_-$. When $\beta+g\mathsf f=0$ the first row of \eqref{eq:agy-identity} is closed, $\E[\bar\mu_{t+u}\mid\text{process up to }t]=e^{-\mathsf p_1u}\bar\mu_t$ with $\mathsf p_1=\varepsilon+g$, and $\varpi_-=\min(\mathsf p_1,\mathsf p_2)$. For the comparison value, $\varpi_+-\mathsf p_1=\frac12[(\mathsf p_2-\mathsf p_1)+((\mathsf p_1-\mathsf p_2)^2+4\mathsf E)^{1/2}]\ge0$ and $\varpi_-\varpi_+=\det\mathsf A=\mathsf p_1\mathsf p_2-\mathsf E$, so
\begin{equation*}
\varpi_-=\frac{\mathsf p_1\mathsf p_2-\mathsf E}{\varpi_+}\le\mathsf p_2-\frac{\mathsf E}{\mathsf p_1}=\delta+\frac{\omega(\varepsilon+g(1-\mathsf f))}{\varepsilon+\beta+g}.
\end{equation*}
When $\mathsf p_2\le\mathsf p_1/2$, the inequality $(1+x)^{1/2}\le1+x/2$ for $x\ge0$, with $x=4\mathsf E/(\mathsf p_1-\mathsf p_2)^2$, gives $\varpi_-\ge\mathsf p_2-\mathsf E/(\mathsf p_1-\mathsf p_2)$. Since $\mathsf E\le\omega\mathsf p_1$ and $\mathsf p_1-\mathsf p_2\ge\mathsf p_1/2$, one has $\mathsf E/(\mathsf p_1-\mathsf p_2)-\mathsf E/\mathsf p_1=\mathsf E\mathsf p_2/(\mathsf p_1(\mathsf p_1-\mathsf p_2))\le2\omega\mathsf p_2/\mathsf p_1$. Hence $\varpi_-$ falls below $\mathsf p_2-\mathsf E/\mathsf p_1$ by at most $2\omega\mathsf p_2/\mathsf p_1$. Without a record, communication among the members conserves their mean and a replacement moves a position from $\mu_i$ to $\mathsf f\bar\mu$ plus mean-zero noise, so the expected change of $\bar\mu$ is $[-\varepsilon\bar\mu-g(1-\mathsf f)\bar\mu]\mathrm dt$.

(b) By \cref{prop:col-records-partition} a record written only by members is internal or boundary, a collective variable whose persistence \cref{def:agy-identity} counts. A record written only by non-members is external. When the writes do not depend on the members' states, the record's path is a function of the exterior's exogenous input $\mathcal X_{\mathrm{ext}}$ of \cref{def:agy-identity}, and $\mathcal X_{\mathrm{ext}}$ is independent of the members' private noises, the replacement events and the newcomers' noises, which are the remaining exogenous inputs of the model. By the first row of \eqref{eq:agy-identity}, $\bar\mu_{t+u}=e^{-(\varepsilon+\beta+g)u}\bar\mu_t$ plus a functional of the record's path plus martingale increments after $t$, which have zero conditional mean given $\mathcal X_{\mathrm{ext}}$ and the past by that independence; so the conditional expectation of $\bar\mu_{t+u}$ given the past and $\mathcal X_{\mathrm{ext}}$ relaxes at the rate $\varepsilon+\beta+g$, and the conditional identity time of $\bar\mu$ is $1/(\varepsilon+\beta+g)\le1/g$.
\end{proof}

\begin{proof}[Proof of \cref{prop:agy-writer-leak}]
With no blanket variables in the description, the leak rate is the mutual information rate of the stationary processes $\mu_{\mathrm{in}}$ and $\mu_{\mathrm{ex}}$. For jointly stationary Gaussian processes with spectral densities bounded and bounded away from zero, this rate is $-\frac12\langle\ln(1-\lvert\mathrm{coh}(\psi)\rvert^2)\rangle_{\mathbb T}$, with $\mathrm{coh}$ the coherence and $\langle\cdot\rangle_{\mathbb T}$ the average over frequencies $\psi$ on the unit circle, a formula attributed to \citet{gelfand1959calculation}. The exterior agent's own filter $(1-c_{\mathrm{ex}}\mathrm L)^{-1}$, with $\mathrm L$ the lag, is invertible and leaves the coherence unchanged, so the rate equals that between $\mu_{\mathrm{in}}$ and $\beta z+\zeta_{\mathrm{ex}}$. With $\mathsf D(\psi)=\lvert1-(1-\delta_z)e^{-\mathrm i\psi}\rvert^2=1+(1-\delta_z)^2-2(1-\delta_z)\cos\psi$, the spectral densities give
\begin{equation*}
1-\lvert\mathrm{coh}(\psi)\rvert^2=\frac{\beta^2q_{\mathrm{rec}}+q_{\mathrm{rd}}\mathsf D(\psi)}{\beta^2(q_{\mathrm{rec}}+\omega^2q_{\mathrm{wr}})+q_{\mathrm{rd}}\mathsf D(\psi)}=\frac{\mathsf a_0-2\mathsf a_\times\cos\psi}{\mathsf a_1-2\mathsf a_\times\cos\psi}.
\end{equation*}
For $\mathsf a\ge2\lvert\mathsf a_\times\rvert$, $\mathsf a-2\mathsf a_\times\cos\psi=\mathsf c_{\mathsf a}\lvert1-\mathsf r_{\mathsf a}e^{\mathrm i\psi}\rvert^2$ with $\lvert\mathsf r_{\mathsf a}\rvert\le1$ and $\mathsf c_{\mathsf a}=(\mathsf a+(\mathsf a^2-4\mathsf a_\times^2)^{1/2})/2$, and $\langle\ln\lvert1-\mathsf r_{\mathsf a}e^{\mathrm i\psi}\rvert^2\rangle_{\mathbb T}=0$, so $\langle\ln(\mathsf a-2\mathsf a_\times\cos\psi)\rangle_{\mathbb T}=\ln\mathsf c_{\mathsf a}$. Both $\mathsf a_0$ and $\mathsf a_1$ exceed $2\lvert\mathsf a_\times\rvert$, since $1+(1-\delta_z)^2\ge2\lvert1-\delta_z\rvert$; this gives \eqref{eq:agy-writer-leak}, which is positive when $\mathsf a_1>\mathsf a_0$, that is, when $\beta\omega\ne0$. With $q_{\mathrm{rec}}=0$, $\mathsf a_0+(\mathsf a_0^2-4\mathsf a_\times^2)^{1/2}=2q_{\mathrm{rd}}$ and $\mathsf a_1^2-4\mathsf a_\times^2=q_{\mathrm{rd}}^2(\delta_z^2+\mathrm{snr})((2-\delta_z)^2+\mathrm{snr})$. Writing $\mathsf s_1=(\delta_z^2+\mathrm{snr})^{1/2}$ and $\mathsf s_2=((2-\delta_z)^2+\mathrm{snr})^{1/2}$, one has $\mathsf a_1/q_{\mathrm{rd}}=(\mathsf s_1^2+\mathsf s_2^2)/2$, so the rate is $\ln((\mathsf s_1+\mathsf s_2)/2)$, the first form of \eqref{eq:agy-writer-regimes}. Since $\mathsf s_1-\delta_z=\mathrm{snr}/(\mathsf s_1+\delta_z)$ and $\mathsf s_2-(2-\delta_z)=\mathrm{snr}/(\mathsf s_2+2-\delta_z)$, $(\mathsf s_1+\mathsf s_2)/2=1+\frac12(\mathsf s_1-\delta_z)(1+(\mathsf s_1+\delta_z)/(\mathsf s_2+2-\delta_z))$; for $\delta_z\le1$ the denominator is at least $2$, so the correction is at most $\mathsf s_1$, and $\ln(1+x)=x(1+O(x))$ with $x=O(\mathsf s_1)$ gives the second form. For $\mathrm{snr}\ll\delta_z^2$, the exact small-$\mathrm{snr}$ limit is $\frac{\mathrm{snr}}4(\frac1{\delta_z}+\frac1{2-\delta_z})=\mathrm{snr}/(2\delta_z(2-\delta_z))$; for $\delta_z^2\ll\mathrm{snr}\ll1$, $\mathsf s_1\approx\mathrm{snr}^{1/2}$.
\end{proof}

\begin{proof}[Proof of \cref{prop:agy-nulls-vanish}]
(a) Given $\mathcal X_{\mathrm{sh}}$, the paths are functions of the private inputs $\zeta_i$, which are independent of each other and of $\mathcal X_{\mathrm{sh}}$, so the paths are conditionally independent; every pairwise term of \eqref{eq:agy-omega} vanishes, and the conditional covariance of the members' states is block diagonal across agents, and so is its inverse.

(b) Because the member reads no message and no record, its path is a function of the world path, the initial state and its own private inputs. Because it neither acts, writes a record nor is read, its private inputs influence no other variable, so they are independent of the world path, the initial state and every other agent's variables. Given the world path and the initial state, its path is therefore conditionally independent of every other variable, and its terms in the statistic conditioned on them vanish.

(c) For a finite window with $X_0$ in the conditioning set, each path $\mu_i[0{:}\bar t]$ is $\sum_kL_{ik}\zeta_k[0{:}\bar t]$ plus a function of $\mathcal X_{\mathrm{sh}}$, with loadings $L_{ik}$ whose entries are polynomials in the coefficients, so the conditional cross-covariance $\Cov(\mu_i[0{:}\bar t],\mu_j[0{:}\bar t]\mid\mathcal X_{\mathrm{sh}})=\sum_kL_{ik}\Cov(\zeta_k)L_{jk}^\top$ has polynomial entries. For Gaussian vectors the conditional mutual information is positive exactly when the conditional cross-covariance is nonzero. Take the directed path $j=k_0\to k_1\to\dots\to k_\ell=i$ in the support, set every coefficient to zero except the path coefficients, which are set to one, and the loading of $\mu_j$ on its private input, which is kept. Then, given $\mathcal X_{\mathrm{sh}}$, $\mu_j(t)$ is a nonzero linear function $\mathsf l^\top\zeta_j(t-1)$ of its private input and $\mu_i(t+\ell)=\mu_j(t)$, so $\Cov(\mu_i(t+\ell),\mu_j(t)\mid\mathcal X_{\mathrm{sh}})=\mathsf l^\top\Cov(\zeta_j(t-1))\mathsf l>0$ by the positive variance of that input. Some entry of the conditional cross-covariance is thus a polynomial that is not identically zero, and its zero set is a proper real algebraic set, closed, of measure zero and with empty interior. For scalar states with independent scalar private inputs and coefficients that are nonnegative and positive on the path and on the loading, every term of $\sum_kL_{ik}\Cov(\zeta_k)L_{jk}^\top$ is nonnegative and the term through $j$'s input carried by the path is positive. For the stationary version, parametrize by the update matrix $W$ on its support with the private noise variances fixed and positive: the window covariances are rational in $W$ on the set $\{\rho(W)<1\}$, solutions of a Lyapunov equation, and this set is open and star-shaped about $0$, since $\rho(sW)=s\rho(W)$ for $s\in[0,1]$, hence connected. A real-analytic function on a connected open set that vanishes on a set of positive measure vanishes identically, and the path matrix above is nilpotent, hence stable, with $\Cov(e_i(t+\ell),e_j(t))=q_j\ne0$, because no edge enters $j$.

(d) With $\Pi=\diag(\mathsf d)$ on a region of $n$ members, $\xout=\bar{\mathsf d}$, the mean of $\mathsf d$ over the region, and $\xin$ is the smallest eigenvalue of the compression of $\diag(\mathsf d)$ to $\one^\perp$, whose trace is $(n-1)\bar{\mathsf d}$; so $\xin\le\bar{\mathsf d}=\xout$ and $r\le1$, with equality when $\mathsf d$ is constant. Under $\mathfrak N_{\mathrm{lp}}$ and $\mathfrak N_{\mathrm{rn}}$ every channel from another agent delivers noise that is either private to the receiver or a sender-side broadcast noise, and under $\mathfrak N_{\mathrm{eff}}$ the world is exogenous. Without records, in the stationary regime of the stable null dynamics each member's state is a measurable function of the world path, the broadcast noises and its own private inputs, the initial state being forgotten. The world-conditional law of \cref{def:col-world-conditional} conditions on the world path and the broadcast noises, so it is the law given $\mathcal X_{\mathrm{sh}}$, and the first part applies.
\end{proof}

\begin{proof}[Proof of \cref{prop:agy-tests}]
(a) Let $U_i$ and $U_j$ be as in the statement and $Z_{\mathrm{sh}}$ the recorded drive, with maximum-likelihood covariance $S$ from the $n_{\mathrm{rep}}$ windows and $S_{\cdot\mid Z}$ the corresponding residual covariances. A coordinate of a path that is a function of the drive has zero residual, so it must be excluded for the residual covariances to be nonsingular; since $\Cov((U_i,U_j)\mid\mathcal X_{\mathrm{sh}})\succ0$ and $n_{\mathrm{rep}}>\dim\mathcal X_{\mathrm{sh}}+\dim U_i+\dim U_j+1$, the residual covariances $S_{U_i\mid Z}$, $S_{U_j\mid Z}$ and $S_{U_iU_j\mid Z}$ are nonsingular with probability one, and the Gaussian regression model of $(U_i,U_j)$ on the drive is regular. Because the excluded coordinates are functions of the drive, $I(U_i;U_j\mid\mathcal X_{\mathrm{sh}})$ equals the conditional mutual information of the full paths. The maximized Gaussian log-likelihood of a regression model is $-\frac{n_{\mathrm{rep}}}2\ln\det$ of its residual covariance plus a constant that depends only on dimensions. Under the hypothesis, the joint density factorizes as $p(z)p(u_i\mid z)p(u_j\mid z)$ and its likelihood is maximized factor by factor, giving $-\frac{n_{\mathrm{rep}}}2[\ln\det S_Z+\ln\det S_{U_i\mid Z}+\ln\det S_{U_j\mid Z}]$; without the hypothesis the maximum is $-\frac{n_{\mathrm{rep}}}2[\ln\det S_Z+\ln\det S_{U_iU_j\mid Z}]$; the constants agree because the dimensions add up. Twice the difference is $n_{\mathrm{rep}}\ln(\det S_{U_i\mid Z}\det S_{U_j\mid Z}/\det S_{U_iU_j\mid Z})=2n_{\mathrm{rep}}\hat I$. The hypothesis sets the $\dim U_i\times\dim U_j$ block of the conditional cross-covariance to zero, a smooth constraint of that codimension in a regular model, so by \citet{wilks1938large} the statistic converges in law to the chi-squared law with $\dim U_i\cdot\dim U_j$ degrees of freedom. Under a causally separating null the hypothesis holds by \cref{prop:agy-nulls-vanish}(a).

(b) Break ties uniformly at random. Under exchangeability the rank of $\mathcal S_0$ among $\mathcal S_0,\dots,\mathcal S_{n_{\mathrm{null}}}$, counted from the top, is uniform on $\{1,\dots,n_{\mathrm{null}}+1\}$, and it is at most $1+\#\{l\ge1:\mathcal S_l\ge\mathcal S_0\}$. The event $p\le\alpha$ says that this count is at most $\alpha(n_{\mathrm{null}}+1)-1$, so it implies that the rank is at most $\lfloor\alpha(n_{\mathrm{null}}+1)\rfloor$, an event of probability $\lfloor\alpha(n_{\mathrm{null}}+1)\rfloor/(n_{\mathrm{null}}+1)\le\alpha$.
\end{proof}

\begin{proof}[Proof of \cref{thm:agy-h1}]
(a) The first sentence is \cref{thm:col-path-blanket}(a) and (b): for the canonical partition with per-channel noise the first two conditions of \eqref{eq:col-pb} hold by construction, the start supplies the fourth, and the third, the collider condition, holds when no sensory channel reads both sides in one round and fails outside a set of coefficients of positive codimension when one does, so the converse holds only generically. The rest is \cref{thm:col-boundary}(a) to (c) with \cref{prop:col-selects-nothing}. (b) The order parameter vanishes under the deletion nulls by \cref{prop:agy-nulls-vanish}(a), and in a population without records no region has contrast above one by \cref{prop:agy-nulls-vanish}(d), so under the nulls no region is selected at a level $r_0>1$. At every parameter of the given open set the full model selects $C$, a collective at the level $r_0$; the selected partitions therefore differ there, for every admissible level in $(1,r(C)]$, which is the clause of \cref{def:agy-nulls}. The positivity of $\Omega_{\mathrm{int}}$ is \cref{prop:agy-nulls-vanish}(c). (c) Apply \cref{thm:agy-closure-average}(a) with $M=\Pi^c_{CC}$, $\xin=r(C)\xout$ and $\sigma_C=s\xout$: then $\xin-\xout=(r(C)-1)\xout>0$, $\sigma_C/(\xin-\xout)=s/(r(C)-1)$ and $\xout/\xin=1/r(C)$, so \eqref{eq:agy-markov-sym} becomes \eqref{eq:agy-h1-closure}; the non-reversible case is \cref{thm:agy-closure-average}(b) with \cref{prop:agy-drift-precision}. (d) is \cref{prop:agy-consensus,prop:agy-literal-consensus}. (e) is \cref{thm:agy-agency,thm:agy-heterogeneous,thm:agy-joint-acts}; for the last sentence, members without data sharing have expected free energies that depend on their own policies only, so the stationary law of deliberation is the product law $\prod_iE_ie^{-\gamma G_i}/\mathsf Z_i$, whose log odds ratios vanish, while by \cref{prop:agy-performative-sign} a group posterior with product habits, a slot-product preference and $\gamma_C>0$ has the negative log odds ratio $-\gamma_CI_{q_C}(o_{A'};o_{B'}\mid o_{A_{-ij}})$ for two members whose slots, opened from the null policies, carry redundant information.
\end{proof}
